%%%%%%%% ICML 2026 Workshop Submission %%%%%%%%%%%%%%%%%%%%%%%%%%%%%
%
% Workshop-track submission using the official ICML 2026 style:
%   main body + unlimited references/appendix, double-blind.
%   The icml2026 style file auto-strips author information unless
%   the [accepted] option is provided. Author names are kept symbolic
%   here for safety.
%
% Style files (icml2026.sty / icml2026.bst / algorithm.sty /
% algorithmic.sty / fancyhdr.sty) live alongside this .tex file.
% Bibliography lives in pivot_ab_references.bib (verified entries).
%
%%%%%%%%%%%%%%%%%%%%%%%%%%%%%%%%%%%%%%%%%%%%%%%%%%%%%%%%%%%%%%%%%%%%%

\documentclass{article}

% Recommended, but optional, packages for figures and better typesetting:
\usepackage{microtype}
\usepackage{graphicx}
\usepackage{subcaption}
\usepackage{booktabs}
\usepackage{multirow}

% hyperref makes hyperlinks in the resulting PDF.
\usepackage{hyperref}

% Use the following line for the initial blind version submitted for review:
\usepackage{icml2026}

% For preprint, use:
% \usepackage[preprint]{icml2026}
%
% If accepted, instead use the following line for the camera-ready submission:
% \usepackage[accepted]{icml2026}

\usepackage{amsmath}
\usepackage{amssymb}
\usepackage{mathtools}

\usepackage[capitalize,noabbrev]{cleveref}

% Running title for ICML page-headers (under 40 chars).
\icmltitlerunning{Causally-Faithful Radiology Report Generation}

\begin{document}

\twocolumn[
  \icmltitle{Causally-Faithful Radiology Report Generation \\
             via Dual-Adversarial Counterfactual Training}

  % Anonymous author block. The icml2026 style file auto-strips this
  % in blind mode (i.e., unless [accepted] is set).
  \begin{icmlauthorlist}
    \icmlauthor{Anonymous Authors}{anon}
  \end{icmlauthorlist}
  \icmlaffiliation{anon}{Anonymous Institution}

  \icmlcorrespondingauthor{Anonymous Authors}{anon@example.org}

  \icmlkeywords{evidence-blindness, radiology report generation,
                causal-faithfulness, dual-adversarial training,
                medical AI safety}

  \vskip 0.3in
]

% This is necessary in blind mode so that the affiliation line is hidden:
\printAffiliationsAndNotice{}

\begin{abstract}
Vision-language models that draft radiology reports often produce text
that does not actually reflect the underlying chest X-ray. Diagnostic
studies have documented this \emph{evidence-blindness}, and existing
fixes wrap the trained model in inference-time filters or conformal
selection without changing what the model has learned. We instead
intervene at training time. We propose
\textbf{dual-adversarial causal-faithfulness training}, a LoRA fine-tune
of MedGemma-4B-IT with two ingredients: a symmetric adversarial filter
that downweights training examples that a text-only or an image-only
predictor can already solve, and a causal-faithfulness loss that
penalises the model whenever its probability of the gold token fails
to drop under image masking. We pair the recipe with
\textbf{IMG\_correct}, a refinement of the standard image-masking-gap
diagnostic that excludes a known failure mode in which a consistency-
loss-trained model merely flips its top prediction under masking
without losing support for the right answer. On the public-data OpenI
chest-X-ray collection, our method drives the image-masking gap from
near zero to roughly $75$\,pp at matched token-level accuracy across
three seeds. A complementary \emph{cross-counterfactual} evaluation,
which we run as an honesty check, finds the gain is largely specific to
the counterfactual the loss trained on: image cropping and centre
occlusion collapse the gap, so the model has tightly fit its training-
time counterfactual rather than the image as a whole. Applying
IMG\_correct retroactively to a set of $18$ previously-reported
verifier configurations reveals that the methodological refinement
\emph{verifies} prior numbers rather than overturning them. We release
training data, code, and weights.
\end{abstract}

\section{Introduction}
\label{sec:intro}

A growing body of work documents that multimodal medical AI systems --- vision-
language models for radiology, ophthalmology, and pathology --- make their
predictions in ways that disregard the visual input. \citet{restrepo2025sms}
introduce Selective Modality Shifting (SMS) on MIMIC-CXR and FairVLMed and
show that six open-source VLMs, including those fine-tuned for medical use,
exhibit a marked dependency on text input that persists despite the
presence of complementary visual information. This phenomenon is the
medical-imaging analogue of the partial-input shortcuts documented for
natural language inference \citep{gururangan2018annotation,
poliak2018hypothesis}, and the multimodal-reward-model shortcuts documented
by \citet{li2025sfc}.

Two facts about the existing literature motivate the present work.
\textbf{First}, evidence-blindness diagnostics in medical imaging have been
\emph{measurement-only}: \citet{restrepo2025sms} explicitly leave training-
time mitigation to future work, and the conformal-selection literature
\citep{gui2024conformal,li2025conflvlm,mohri2024conformal} similarly
wraps an existing predictor without modifying the underlying model.
\textbf{Second}, the only training-time mitigation applied at scale to
multimodal models --- the Shortcut-Failure Coefficient (SFC) reweighting of
\citet{li2025sfc} --- operates only on the text-only side, despite the
fact that an image-only shortcut is just as conceptually possible and, on
naturalistic medical data, just as likely. To our knowledge, no published
work combines a text-only and image-only adversarial filter for a medical
multimodal generator and demonstrates an effect on a counterfactual-
grounding diagnostic.

Our \textbf{two primary contributions} are as follows. \textbf{First}, we
develop and release a fine-tuning recipe --- \emph{dual-adversarial causal-
faithfulness training} --- that addresses both training-distribution
shortcuts symmetrically and adds an online causal-faithfulness loss with a
token-wise correctness-gain margin. The full method is a LoRA fine-tune of
MedGemma-4B-IT \citep{google2024medgemma} trained on OpenI
\citep{demner2016openi} using only public-data sources.
\textbf{Second}, we contribute a methodological patch to the standard
image-masking gap (IMG): IMG\_correct, which counts only those examples
where the model's correct-class probability dropped by a margin $\tau$
under image masking, ruling out the \emph{induced inversion} failure mode
in which a consistency-loss-trained model is wrong on the image-zeroed input
by flipping to a different-but-still-wrong answer rather than dropping its
support for the gold class. We apply IMG\_correct retroactively to our
prior verifier work and find that, contra the pre-registered concern, the
induced-inversion contribution is uniformly small. As a \emph{third
instrumentation release} --- not a validated contribution --- we ship a
training-time gradient-norm monitor; we explicitly do not claim a deployable
monitor in the present paper because the held-out mech-interp sweep needed
to validate its predictive AUROC was beyond our compute envelope.

A \emph{cross-counterfactual} eval (\S\ref{sec:cross-cf}) honestly reports
that the headline $+72.58$\,pp IMG\_correct improvement does not generalise
cleanly to localised counterfactuals (image-cropping, centre-occlusion);
the model has tightly fit the trained per-channel-mean and noise
distributions but does not transfer to counterfactuals it has not seen.
We frame the contribution accordingly: this paper demonstrates that
\emph{training-time counterfactual interventions can be tightly-fit on the
in-distribution counterfactual}, not that the resulting model is
image-grounded in a deeper sense. A more diverse training-time counterfactual
mix is the natural follow-up.

\section{Related Work}
\label{sec:related}

\paragraph{Evidence-blindness diagnostics.} \citet{restrepo2025sms}
introduce the Selective Modality Shifting (SMS) family of perturbation-
based diagnostics that swap images or text between samples with opposing
labels. Our IMG\_correct is conceptually a refinement of the
image-perturbation-only branch of SMS, with an additional gold-class-
probability-drop margin that excludes induced inversion. Earlier diagnostic
work for natural-language entailment partial-input baselines
\citep{gururangan2018annotation, poliak2018hypothesis} provides the
conceptual precedent for treating shortcut-solvability as a property of
the dataset rather than the model.

\paragraph{Adversarial filtering for shortcuts.} \citet{li2025sfc}
introduce the Shortcut-Failure Coefficient, a continuous training-sample
reweighting scheme for multimodal reward models that down-weights samples
solvable by a text-only model. The SFC is operated only along the text
axis; the obvious complementary mechanism --- an image-only shortcut
detector that mirrors the text-only SFC --- is not implemented in
\citet{li2025sfc} or, to our knowledge, any subsequent work for the
medical-imaging setting. We extend the family with a \emph{symmetric dual
filter} aggregated as $\min(w_{\text{text}}, w_{\text{image}})$ to
downweight rows that are unimodally solvable along \emph{either} axis.

\paragraph{Conformal selection on top of frozen multimodal models.}
\citet{gui2024conformal} train an alignment predictor on a reference
dataset of human-labelled alignment scores and apply a data-dependent
threshold at deployment. \citet{li2025conflvlm} treat each generated
text claim as an individual hypothesis and demonstrate, on LLaVa-1.5
across scene/medical/document settings, that the error rate can be reduced
from $87.8\%$ to $10.0\%$ at a $95.3\%$ true-positive rate. Both methods
operate on a frozen underlying model. Our contribution is orthogonal: we
improve the underlying generator's faithfulness, which compounds with
rather than competes with any subsequent conformal wrapper.

\paragraph{Audit-and-repair frameworks for radiology generators.} RadFlag
\citep{zhang2025radflag} samples reports at varying temperatures and uses
self-agreement to flag hallucinated claims; ReXTrust
\citep{hardy2025rextrust} trains a self-attention head over LVLM hidden
states (AUROC $0.875$); FactCheXcker \citep{heiman2025factchexcker}
de-hallucinates measurement statements via a code-generation pipeline.
All three are \emph{post-hoc, inference-time} hallucination filters
operating on a frozen generator; none implements a training-time
intervention against evidence-blindness. The training-time, model-
modifying contribution proposed here stacks with each (a more faithful
underlying generator delivers fewer hallucinations to filter regardless
of the filtering recipe).

\paragraph{Methodological-audit precedent.} \citet{recht2019imagenet}
re-audit ImageNet-trained classifiers under a held-out-from-the-original-
distribution test set and find that aggregate accuracy claims do not
transfer. Our IMG\_correct audit of prior IMG numbers is the medical-
imaging analogue: we accept that prior IMG numbers are correct under the
IMG definition but ask how much of the reported gap survives under a
more demanding definition that excludes induced inversion.

\section{Methods}
\label{sec:methods}

\subsection{Backbone and parameter-efficient fine-tuning}
\label{sec:backbone}

We fine-tune MedGemma-4B-IT \citep{google2024medgemma}, a 4.31\,B-parameter
image-text instruction-tuned variant of Gemma-3 trained on a mixture of
medical images including chest radiographs, pathology slides, dermatology
images, and fundus photography. The model exposes a standard
\texttt{AutoModelForImageTextToText} interface with a \texttt{pixel\_values}
keyword argument that accepts a $(B, 3, H, W)$ tensor of normalised image
patches. We wrap the base model with a LoRA adapter \citep{hu2022lora} of
rank $16$ and alpha $32$, applied to the
\texttt{q\_proj}/\texttt{k\_proj}/\texttt{v\_proj}/\texttt{o\_proj} modules
of every transformer block. Trainable parameter count: $11.9$\,M ($0.28\%$
of base); the rest is frozen. CheXagent-2 \citep{chen2024chexagent} was the
original target backbone but its tokenizer-embedded image paradigm does
not accept a \texttt{pixel\_values} keyword, blocking counterfactual
training without substantial reverse-engineering effort. MAIRA-2
\citep{bannur2024maira2} requires similar backbone-specific work that we
did not undertake. Our backbone choice was driven by this engineering
constraint.

\subsection{Composite training loss}
\label{sec:loss}

The training objective combines three components,
\begin{equation}
\mathcal{L}_{\mathrm{total}}(x) = w(x)\cdot\mathcal{L}_{\mathrm{SFT}}(x) +
\lambda_{\mathrm{faith}}\cdot\mathcal{L}_{\mathrm{faith}}(x),
\end{equation}
where $w(x)\in[w_{\min},1]$ is the per-row weight produced by the dual
adversarial filter (\S\ref{sec:dual}), $\mathcal{L}_{\mathrm{SFT}}$ is
supervised fine-tuning cross-entropy on the gold radiology-report tokens,
and $\mathcal{L}_{\mathrm{faith}}$ is the causal-faithfulness loss
(\S\ref{sec:faith}). The faithfulness weight $\lambda_{\mathrm{faith}}$ is
a fixed hyperparameter set to $0.3$ throughout, chosen by a small smoke-
search over $\{0.1, 0.3, 1.0\}$ on a 6-step pre-flight.

\subsection{Dual adversarial filter}
\label{sec:dual}

The dual filter combines two independently-trained shortcut detectors: a
text-only branch that scores how predictable the gold label is from the
claim and evidence text alone, and an image-only branch that scores how
predictable the gold label is from the image alone. Each branch produces
a per-row weight in $\{w_{\min}, 1.0\}$ (discrete schedule) or
$[w_{\min}, 1.0]$ (continuous SFC ablation). The dual weight is
$w(x) = \min(w_{\text{text}}(x), w_{\text{image}}(x))$: a row is
downweighted aggressively if it is unimodally solvable along \emph{either}
axis.

The text-only branch reuses the implementation from our prior verifier
work [citation withheld for double-blind review]: a RoBERTa-large
cross-encoder trained for one epoch on (claim, evidence)\,$\to$\,gold-
label pairs from the training split. Rows where the text-only model's
predicted probability of the gold class exceeds $0.7$ are downweighted
to $w_{\min} = 0.2$; the remaining rows are kept at $1.0$.

The image-only branch is the genuine novelty axis. We extract image
features via the BiomedCLIP image encoder \citep{zhang2024biomedclip} and
train a 2-layer MLP head on image\,$\to$\,gold-label for one epoch. The
same threshold rule is applied. On the $20{,}560$-row training pool, the
text-only branch downweights $85.3\%$ of rows; the image-only branch
$81.2\%$; the union after $\min$ aggregation $91.7\%$. Crucially,
$\mathbf{6.5\%}$ of rows are caught by the image-only branch but missed
by the text-only branch (Table~\ref{tab:filter-decomp}, App.~\ref{app:filter})
--- exactly the rows whose downweighting cannot be reproduced by a
text-only filter alone, and the genuine novelty axis relative to
\citet{li2025sfc}. In an ablation we additionally compare the discrete
$\{w_{\min},1\}$ rule against the continuous SFC formulation
$\mathrm{SFC}(x) = 1 - p_{\text{text}}^{\text{correct}}(x)$ clipped to
$[w_{\min}, 1]$.

\subsection{Causal-faithfulness loss with token-wise gain margin}
\label{sec:faith}

The naive consistency loss
$\mathrm{ReLU}(\text{margin} - \mathrm{KL}(p_{\mathrm{full}} \| p_{\mathrm{alt}}))$
admits a degenerate optimum: the model can drive $p_{\mathrm{alt}}$ to an
arbitrary distribution while $p_{\mathrm{full}}$ is shaped by the SFT loss
alone, and the divergence margin can be satisfied without the model ever
using the image. This is the \emph{induced inversion} failure mode that
IMG\_correct directly measures (\S\ref{sec:imgcorrect}). To rule it out at
training time, we replace the symmetric KL with a \emph{correctness-gain
margin} on gold tokens:
\begin{equation}
\mathcal{L}_{\mathrm{faith}}(x) =
\frac{1}{|T|}\!\sum_{t \in T}\mathrm{ReLU}\!\Bigl(
\text{margin} - \bigl(\log p_{\mathrm{full}}(y_t) - \log p_{\mathrm{alt}}(y_t)\bigr)\!\Bigr)
\end{equation}
where $y_t$ is the gold token at position $t$, $p_{\mathrm{full}}$ and
$p_{\mathrm{alt}}$ are the per-position distributions under the full and
image-zeroed inputs, $T$ is the set of assistant-turn token positions, and
$\text{margin} = 0.3$\,nats (a deliberately small value; faithfulness-via-
reduction-in-gold-token-probability is what we want, not faithfulness-via-
arbitrary-distributional-divergence). The token-wise margin couples the
divergence requirement to gold-token correctness on the full input, so the
model cannot satisfy the margin by being arbitrarily wrong on the image-
zeroed input. The same margin acts symmetrically across both image-zeroed
and image-flipped counterfactuals; on rows that are not laterality-
sensitive, the image-flipped counterfactual is replaced with the image-
zeroed counterfactual as a no-op fallback (masked out via
\texttt{flipped\_row\_mask}).

Image-zeroing is implemented as the per-channel mean of the input batch's
pixel values, \emph{not} a literal zero tensor: a literal zero is, after
the processor's normalisation, a specific out-of-distribution input that
the model has never seen at training time. The faithfulness loss requires
three forward passes per training step (full, image-zeroed, image-flipped
on laterality rows). To stay memory-bounded on a single H100 we use
batch size $2$ with gradient-accumulation factor $8$, yielding effective
batch size $16$.

\subsection{IMG\_correct: methodological refinement of the image-masking gap}
\label{sec:imgcorrect}

The standard image-masking gap is
$\mathrm{IMG}(f, \mathcal{D}) = \mathrm{acc}(f, \mathcal{D}) - \mathrm{acc}(f, \mathcal{D}^{\mathrm{img}=0})$.
This counts every example on which the model is correct on the full input
and incorrect on the image-zeroed input as evidence of image-grounded
prediction. But a consistency-loss-trained model can lower its image-zeroed
correctness to zero by always flipping its argmax under masking,
\emph{even if its gold-class probability barely drops}. We call the cases
where the model flips its argmax under masking but does not drop its
gold-class probability by $\geq\tau$ the \emph{induced inversion}
sub-population: observationally identical to faithful grounding under IMG,
but mechanistically distinct.

IMG\_correct addresses this by counting only those examples where (a) the
model is correct on the full input ($\arg\max_{\mathrm{full}} = y$);
(b) the model is incorrect on the image-zeroed input
($\arg\max_{\mathrm{masked}} \neq y$); \emph{and} (c) the gold-class
probability under masking is lower than under the full input by at least
margin $\tau$:
$p_{\mathrm{full}}(y) - p_{\mathrm{masked}}(y) \geq \tau$. We additionally
report \emph{induced inversion} as
$(n_{\text{ci,no-margin}} - n_{\text{ci,with-margin}})/n \times 100$,
isolating the sub-population of ``model flipped under masking but its
correct-class probability did not drop by $\geq\tau$.'' Note IMG\_correct
\emph{can} exceed IMG when many examples transition incorrect-to-correct
under masking; this is mathematically why several configurations in our
retroactive audit (\S\ref{sec:audit}) show IMG\_correct $>$ IMG. For
autoregressive generators we apply the same logic per token on the gold
target sequence and average over assistant-turn positions only.

\section{Experimental Setup}
\label{sec:setup}

\paragraph{Data.} All training and evaluation use public-data sources. Training distribution: OpenI
\citep{demner2016openi} --- $10{,}839$ (claim, image, gold-report) tuples
after our text-target filter. In-distribution evaluation: standard OpenI
test split ($1{,}587$ examples). Cross-site eval: ChestX-Det10 silver-
target subset ($1{,}122$ rows) \citep{liu2020chestxdet}; PadChest-GR
\citep{castro2024padchestgr} is deferred to follow-up because no silver
targets are available. We do not use MIMIC-CXR \citep{johnson2019mimiccxr}
in this study.

\paragraph{Configurations.} Five configurations $\times$ three seeds:
\textbf{sft\_only} (SFT-only baseline); \textbf{sft\_dual} (SFT $+$ dual
filter, no faith loss); \textbf{sft\_faith} (SFT $+$ faith loss, no dual
filter); \textbf{sft\_full} (full method: SFT $+$ faith $+$ dual filter,
the headline configuration); \textbf{sft\_full\_sfc} (full method with the
continuous \citet{li2025sfc}-style SFC reweighting). Each runs at seeds
$\{42, 1337, 9001\}$. Total training cost $\approx \$200$ on Modal H100
80\,GB.

\paragraph{Metrics.} Primary: \textbf{IMG\_correct} at margin $\tau = 0.1$.
Secondary: standard IMG, induced-inversion contribution, accuracy on
supervised target tokens. We do not report report-quality metrics (BLEU,
ROUGE, BERTScore-style overlap) as primary because the OpenI-only training
distribution makes those a single-site comparison rather than a
generalisation claim; the central methodological claim is the cross-
counterfactual diagnostic.

\section{Results}
\label{sec:results}

\subsection{IMG\_correct on the headline five-config ablation}
\label{sec:headline}

We report IMG\_correct, IMG, induced inversion, and accuracy for each of
the $5\times3 = 15$ Pivot fine-tunes against the SFT-only seed-averaged
baseline. The pre-registered prediction is that the full-method
configuration achieves IMG\_correct strictly greater than SFT-only by at
least $5$\,pp at matched accuracy. All $15$ fine-tunes have completed.
Numbers in Table~\ref{tab:headline} are seed-means and standard deviations
over the three independent seeds per configuration, aggregated over
$33{,}173$ valid assistant-turn token positions across the OpenI test
split ($\approx 21$ tokens/example).

\begin{table}[t]
\centering
\small
\caption{\textbf{Headline IMG\_correct on the $1{,}587$-example OpenI test
split (token-level, $n=33{,}173$).} Seed-mean $\pm$ standard deviation
over three seeds (42, 1337, 9001) per configuration. Induced inversion is
essentially zero across all faith-loss configurations.}
\label{tab:headline}
\setlength{\tabcolsep}{2.6pt}
\begin{tabular}{lcccc}
\toprule
\textbf{Config} & \textbf{IMG\_corr (pp)} & \textbf{IMG (pp)} & \textbf{Inv (pp)} & \textbf{Acc (\%)}\\
\midrule
sft\_only         & $2.38\pm0.26$  & 0.36  & 0.40 & 74.31 \\
sft\_dual         & $2.10\pm0.24$  & 0.89  & 0.58 & 75.41 \\
\textbf{sft\_faith}    & $\mathbf{74.11\pm0.26}$ & 74.04 & 0.01 & 75.61 \\
\textbf{sft\_full}     & $\mathbf{74.96\pm0.83}$ & 74.96 & 0.01 & 75.48 \\
\textbf{sft\_full\_sfc}& $\mathbf{75.42\pm0.28}$ & 75.43 & 0.01 & 75.44 \\
\bottomrule
\end{tabular}
\end{table}

\textbf{The pre-registered $\geq 5$\,pp threshold is exceeded by an order
of magnitude on three independent configurations $\times$ three independent
seeds each.} The full method gives a $\mathbf{+72.58}$\,\textbf{pp}
\textbf{improvement at $+1.17$\,pp accuracy} over the SFT-only baseline; the
continuous-SFC variant gives $+73.04$\,pp at $+1.13$\,pp; the faith-loss-
alone configuration gives $+71.73$\,pp at $+1.30$\,pp. Reading down the
table, the dual filter alone (sft\_dual) is statistically indistinguishable
from baseline; the faith loss (sft\_faith) is the load-bearing component;
adding the dual filter on top of the faith loss contributes a further
$+0.85$\,pp seed-mean (within one standard deviation of sft\_full); the
continuous-SFC variant contributes a further $+0.46$\,pp on top of that.
Induced-inversion contributions are essentially zero across all faith-
loss configurations ($\le 0.01$\,pp), meaning the entire IMG gap is
faithful grounding under our refined definition.

\subsection{Cross-counterfactual eval}
\label{sec:cross-cf}

The principal limitation of \S\ref{sec:headline} is that the faith loss and
the IMG\_correct metric measure overlapping per-token quantities: a model
that minimises the faith loss is mechanically pushed toward higher
IMG\_correct on the \emph{exact counterfactual distribution} the loss
trained on (per-channel mean image-zeroing). To disambiguate ``the model
learned its training objective'' from ``the model genuinely uses the
image,'' we evaluate the headline sft\_full configuration under three
additional counterfactuals the faith loss did \emph{not} see during
training: (i) \textbf{noise} replaces pixels with Gaussian noise at the
per-channel mean $\pm$ std; (ii) \textbf{crop} keeps only a $30\% \times
30\%$ upper-left patch of the image ($\approx 9\%$ of image area) and
replaces the rest with the per-channel mean; (iii) \textbf{occlude}
zeroes a centred $50\% \times 50\%$ patch ($\approx 25\%$ of image area)
and preserves the periphery.

\begin{table}[t]
\centering
\small
\caption{\textbf{Cross-counterfactual eval on sft\_full (seed 42 only).}
Single-seed numbers; the trained-mean cell ($74.00$\,pp) is one realisation
of the IMG\_correct $= 74.96 \pm 0.83$\,pp distribution from
Table~\ref{tab:headline}. Baseline anchor (sft\_only seed 42 under crop):
IMG\_correct $= 2.80$\,pp at $73.82\%$ accuracy.}
\label{tab:crosscf}
\setlength{\tabcolsep}{4.0pt}
\begin{tabular}{lccc}
\toprule
\textbf{Counterfactual} & \textbf{IMG (pp)} & \textbf{IMG\_corr (pp)} & \textbf{Inv (pp)}\\
\midrule
\textbf{mean (TRAINED)} & 73.97 & 74.00 & 0.01 \\
noise                   & 62.10 & 62.20 & 0.10 \\
crop                    & \phantom{0}3.64 & \phantom{0}4.81 & 0.61 \\
occlude                 & \phantom{0}0.28 & \phantom{0}1.11 & 0.48 \\
\bottomrule
\end{tabular}
\end{table}

The result is a clean partial-circularity confirmation. Under the trained
counterfactual (mean), IMG\_correct is $74.00$\,pp on this seed (consistent
with the 3-seed reading). Under a \emph{similar} untrained counterfactual
(noise: also a globally-uniform-OOD pattern), IMG\_correct stays high at
$62.20$\,pp. Under \emph{dissimilar} untrained counterfactuals (crop, which
preserves $9\%$ of true image; occlude, which zeroes the centre),
IMG\_correct collapses to $4.81$\,pp and $1.11$\,pp respectively. The
matched-counterfactual sft\_only baseline yields IMG\_correct $= 2.80$\,pp
on crop. The cross-counterfactual sft\_full $-$ sft\_only IMG\_correct gap
on crop is therefore $4.81 - 2.80 = 2.01$\,pp, an order of magnitude
smaller than the $72.58$\,pp gap on the trained-mean counterfactual.

The honest reading: \textbf{the $+72.58$\,pp headline IMG\_correct
improvement does not generalise cleanly across counterfactuals}. The model
has demonstrably learned to drive the gold-token probability to $\approx 0$
under the per-channel-mean and noise counterfactuals (which are global,
uniform-OOD distributions), but it has \emph{not} learned a more general
``use the image'' behaviour that transfers to localised counterfactuals.

\subsection{Cross-site transfer (ChestX-Det10 silver)}
\label{sec:crosssite}

The OpenI-only training distribution means the headline numbers are
in-distribution. To test cross-site transfer we evaluate sft\_full and
sft\_only on a ChestX-Det10 \citep{liu2020chestxdet} silver-target subset
generated under our prior radiology-report-generation infrastructure (the
silver labels are CheXagent-derived, not radiologist gold; the eval is a
transfer test, not a faithfulness-against-radiologist test).

\begin{table}[t]
\centering
\footnotesize
\caption{\textbf{Cross-site eval on ChestX-Det10 silver targets.} The
sft\_full cross-site run was capped at $500$ rows; the baseline used the
full $1{,}122$-row silver-target subset. IMG\_correct is a per-token rate
and is not biased by the absolute token count.}
\label{tab:crosssite}
\setlength{\tabcolsep}{2.0pt}
\begin{tabular}{lcccc}
\toprule
\textbf{Config / Site} & \textbf{n\_tok} & \textbf{IMG\_corr} & \textbf{Inv} & \textbf{Acc}\\
\midrule
sft\_full / OpenI (3-seed) & $99{,}519$ & $74.96$ & $0.01$ & $75.48$\\
sft\_full / chestx (silver, $s{=}42$) & $15{,}955$ & $54.62$ & $0.03$ & $54.65$\\
sft\_only / chestx (silver, $s{=}42$) & $35{,}944$ & $\phantom{0}5.12$ & $0.58$ & $50.33$\\
\bottomrule
\end{tabular}
\end{table}

The cross-site sft\_full shows IMG\_correct $= 54.62$\,pp at $54.65\%$
accuracy ($\approx 21$\,pp drop in both metrics relative to in-distribution).
The drop is roughly proportional. The cross-site sft\_only baseline yields
IMG\_correct $= 5.12$\,pp at $50.33\%$ accuracy: it preserves accuracy
within $4$\,pp of in-distribution but barely uses the image. The cross-
site sft\_full $-$ sft\_only IMG\_correct gap is $\mathbf{49.5}$\,\textbf{pp},
about $68\%$ of the in-distribution gap of $72.58$\,pp, consistent with
\textbf{partial transfer rather than collapse}.

\subsection{IMG\_correct retroactive audit on prior systems}
\label{sec:audit}

We apply IMG\_correct to the v5 / v6 verifier checkpoints of our prior
work [citation withheld for double-blind review] (nine verifier
configurations and nine v5 ablations, each on $n=2{,}974$ verifier rows).
Table~\ref{tab:audit} summarises the headline cells; the full audit set
(18 cells) appears in Appendix~\ref{app:audit}.

\begin{table}[t]
\centering
\footnotesize
\caption{\textbf{IMG\_correct retroactive audit (selected cells).} Full
table at Appendix~\ref{app:audit}. The pre-registered concern was that
$\geq 30\%$ of $\mathrm{v5\_3\_contrast}$'s IMG would resolve to induced
inversion; the empirical induced inversion is $1.1\%$ of IMG.}
\label{tab:audit}
\setlength{\tabcolsep}{3.5pt}
\begin{tabular}{lcccc}
\toprule
\textbf{Config} & \textbf{IMG\_corr} & \textbf{IMG} & \textbf{Inv} & \textbf{Inv/IMG}\\
\midrule
v5\_0\_base               & 3.83  & 1.28  & 0.17 & 13.2\%\\
\textbf{v5\_3\_contrast}  & \textbf{71.05} & \textbf{68.12} & \textbf{0.74} & \textbf{1.1\%}\\
v5\_4\_final              & 68.43 & 63.55 & 0.40 & 0.6\%\\
v6\_0\_3site              & 74.51 & 70.21 & 0.07 & 0.1\%\\
\addlinespace
abl\_no\_consist          & 3.60  & 2.79  & 0.30 & 10.8\%\\
abl\_no\_contrast         & 63.28 & 60.02 & 0.87 & 1.4\%\\
abl\_hothresh\_80         & 77.98 & 74.78 & 0.74 & 1.0\%\\
\bottomrule
\end{tabular}
\end{table}

\textbf{Headline finding (negative result; reported transparently).} Across
the nine headline configurations and nine v5 ablations, induced-inversion
contributions are uniformly small: between $0.00$\,pp (v6\_0\_loo\_no\_openi)
and $1.31$\,pp (v6\_0\_loo\_no\_padchest); in relative terms, between
$0.0\%$ and $13.2\%$ of the corresponding IMG. The decisive cell ---
$\mathrm{v5\_3\_contrast}$, the configuration corresponding to the
previously-reported headline IMG $\approx 69$\,pp claim --- yields IMG
$= 68.12$\,pp, IMG\_correct $= 71.05$\,pp, induced inversion $= 0.74$\,pp
(only $1.1\%$ of IMG). The newly-completed $\mathrm{v6\_0\_3site}$ cell ---
the union-of-three-sites configuration --- yields the cleanest reading at
full-data scale: induced inversion $= 0.07$\,pp $= 0.1\%$ of IMG. The
$\mathrm{abl\_no\_consist}$ ablation (consistency loss \emph{ablated
entirely}) is the second-most decisive cell: removing the consistency loss
collapses IMG\_correct to $\approx$ baseline (3.60\,pp) while
\emph{simultaneously} reducing induced inversion (0.30\,pp), the opposite
of what the pre-registered concern (consistency loss \emph{causes} induced
inversion) would predict. \textbf{The pre-registered $\geq 30\%$ hypothesis
is decisively refuted across the entire audited set.}

\section{Discussion}
\label{sec:discuss}

\paragraph{What IMG\_correct reveals about prior systems.} The honest
reading of \S\ref{sec:audit} is that \textbf{the IMG metric, as previously
reported on the v5/v6 verifier family, was already largely measuring real
image-grounded behaviour}. The methodological concern that motivated
IMG\_correct --- that consistency-loss training could in principle reward
the model for becoming wrong-under-masking without using the image ---
remains a real concern \emph{theoretically}, but the empirical evidence
does not support the claim that consistency-loss training produced
substantial induced inversion at the calibration tested. The strongest
version of the IMG\_correct contribution, given this result, is
\emph{not} ``audit reveals overstated grounding claims''; it is (a) a
\emph{verification} that prior IMG numbers held up under a strict
per-example correctness criterion, increasing confidence in prior claims;
and (b) a default-to-report-alongside-IMG diagnostic that gives future
workers a way to detect the failure mode if a less calibrated training
recipe were to introduce it. We avoid the temptation to relax the pre-
registered hypothesis post-hoc.

\paragraph{Cross-counterfactual partial generalisation.} \S\ref{sec:cross-cf}
gives the operational meaning of the circularity caveat: the headline
$+72.58$\,pp number measures responsiveness to the trained counterfactual,
with strong but not complete transfer to similar untrained counterfactuals.
This paper demonstrates that \emph{training-time counterfactual
interventions can be tightly-fit on the in-distribution counterfactual},
not that the resulting model is image-grounded in a deeper sense. A more
diverse training-time counterfactual mix (random crop / occlude / noise /
mean per training step) is the natural follow-up; we expect it would close
the gap.

\paragraph{Cross-site partial transfer.} \S\ref{sec:crosssite} shows
sft\_full's image-grounded behaviour transfers cross-site
(IMG\_correct$=54.62$\,pp on ChestX-Det10 silver, $\approx 68\%$ of in-
distribution gap retained) when the \emph{counterfactual} is held constant
(mean-masking) but the test distribution changes. Cross-counterfactual
generalisation is the dominant remaining failure mode; cross-site
distribution shift is not.

\paragraph{Filter calibration is task-mismatched.} The text-only and
image-only adversarial filters were trained for the verifier task in our
prior work [citation withheld for double-blind review] and reused as
shortcut detectors for the present generator task. The conceptually-
aligned filter for the new task would ask whether a text-only or
image-only generator could reproduce the gold report; our existing filter
is a directional proxy. The dual-filter ablation in
Table~\ref{tab:headline} still demonstrates the mechanism, but the
calibration is approximate.

\paragraph{Public-data-only training is a feature, not a bug.} The
reproducibility literature in medical imaging consistently flags
access barriers as a constraint on independent replication of
MIMIC-CXR-trained results. Our public-data-only commitment means any
researcher with a Hugging Face account and a Modal-or-equivalent compute
provider can re-run the entire pipeline. This restriction costs us scale
(MIMIC-CXR has $\approx 50\times$ more rows than OpenI) but the cost is
recoverable: the methodology generalises to MIMIC-CXR for any researcher
with the relevant access, and the results we report should be read as a
\emph{lower bound} on what is achievable with comparable methods on
larger training distributions.

\section{Limitations}
\label{sec:limits}

\textbf{(L1)} The faith loss (\S\ref{sec:faith}) and the IMG\_correct
metric (\S\ref{sec:imgcorrect}) measure overlapping per-token quantities;
the headline $+72.58$\,pp number is best read as evidence that the model
has fit the training objective on held-out data, not as independent
evidence of image-grounded prediction. \S\ref{sec:cross-cf}
quantifies this caveat empirically. \textbf{(L2)} OpenI-only training;
cross-site transfer is measured on ChestX-Det10 silver targets only.
PadChest-GR cross-site is deferred. \textbf{(L3)} Filter calibration is
task-mismatched (\S\ref{sec:discuss}). \textbf{(L4)} LoRA rank $16$ ($11.9$\,M
trainable parameters of $4.31$\,B); we have not run a rank ablation.
\textbf{(L5)} Single backbone (MedGemma-4B-IT). \textbf{(L6)} Image-masking
is one of several plausible counterfactuals; image-cropping or occlusion
test different aspects of image-grounded prediction and yield different
IMG\_correct numbers. We do not claim image-zeroing is the canonical
counterfactual; we claim that any counterfactual benefits from the
IMG\_correct refinement over the IMG metric. \textbf{(L7)} Cross-system
retroactive audit is on our own prior verifier checkpoints only; closed-
weight commercial vision-language systems and tokenizer-embedded systems
(MAIRA-2, CheXagent-2-3B) are deferred because the per-token softmax
extraction required for IMG\_correct is non-trivial or unavailable under
those paradigms. \textbf{(L8)} Single radiology
modality (chest radiographs); the methodological contribution generalises
to any image-conditioned generation task where text-only and image-only
shortcuts are plausible, but we have not tested this.

\section{Conclusion}
\label{sec:conclusion}

We presented dual-adversarial causal-faithfulness training, a LoRA-based
fine-tuning recipe for image-conditioned medical generators that addresses
evidence-blindness symmetrically along both modalities and rules out the
induced-inversion failure mode at training time via a token-wise
correctness-gain margin loss. We additionally introduced IMG\_correct, a
methodological refinement of the standard image-masking gap. The pre-
registered empirical claims, tested against actual results: \textbf{(i)}
Full method reduces IMG\_correct on OpenI-test by $\geq 5$\,pp at matched
accuracy: \textbf{HIT BY AN ORDER OF MAGNITUDE} on three seeds $\times$
three faith-loss configs ($+72.58$\,pp); \textbf{(ii)} IMG\_correct
refinement reveals $\geq 30\%$ of v5\_3\_contrast's IMG as induced
inversion: \textbf{REFUTED} ($1.1\%$); \textbf{(iii)} dual filter
dominates the faith loss in the ablation: \textbf{INVERTED} (the faith
loss is the load-bearing component). One pre-registration HIT, one
REFUTED, one INVERTED, all reported transparently. The methodological
refinement (IMG\_correct) lands as a contribution independently of the
empirical results, since it is a reusable audit-tool that improves the
interpretation of any future or prior IMG-style metric. We advocate that
counterfactual evidence-sensitivity testing precede aggregate-accuracy
reporting for any multimodal medical generator, and that the IMG\_correct
margin become a default reporting requirement.

\bibliographystyle{icml2026}
\bibliography{pivot_ab_references}

\appendix

\section{Reproducibility Checklist}
\label{app:repro}

\begin{itemize}
\item \textbf{[\checkmark]} Training data sources are public (OpenI for training and in-distribution evaluation; ChestX-Det10 silver-target for cross-site eval; PadChest-GR named as deferred cross-site target; MIMIC-CXR is not used in this study).
\item \textbf{[\checkmark]} Hyperparameters specified per-config in YAML; full hyperparameter table in \S\ref{app:hyper}.
\item \textbf{[\checkmark]} 30 unit tests pass on CPU; full test list in \S\ref{app:tests}.
\item \textbf{[\checkmark]} Modal infrastructure described; all entrypoints documented.
\item \textbf{[\checkmark]} Per-fine-tune training-step counts, save-every-steps schedule, and resume-from-checkpoint logic released as code.
\item \textbf{[\checkmark]} Cross-site IMG\_correct on ChestX-Det10 silver-target measured (\S\ref{sec:crosssite}).
\item \textbf{[\,\,]} Code release URL --- anonymized for double-blind review; will be filled at camera-ready.
\item \textbf{[\,\,]} Trained model weights release URL --- anonymized for double-blind review; subject to MedGemma research-only redistribution terms.
\item \textbf{[\,\,]} Cross-site IMG\_correct on PadChest-GR (gold or silver) --- see \S\ref{sec:limits}; deferred to follow-up.
\item \textbf{[\,\,]} Three-seed averaging on noise/crop/occlude cross-counterfactuals --- single-seed in \S\ref{sec:cross-cf}; deferred to follow-up.
\item \textbf{[\,\,]} Mech-interp sweep validating the gradient-norm monitor's AUROC --- deferred to follow-up.
\end{itemize}

\section{Extended Methods}
\label{app:methods}

\subsection{Faith loss derivation, edge cases, and implementation}
\label{app:faith}

The naive consistency loss is
$\mathcal{L}_{\mathrm{KL}} = \mathrm{ReLU}(\text{margin} - \mathrm{KL}(p_{\mathrm{full}} \| p_{\mathrm{alt}}))$.
This admits a degenerate optimum: a model that always emits a uniform
distribution under masking satisfies any finite KL margin without using
the image. The token-wise correctness-gain margin replaces the
distributional divergence with a per-token signal:
$\mathcal{L}_{\mathrm{faith}}(x) = (1/|T|)\sum_{t \in T}\mathrm{ReLU}(\text{margin} - (\log p_{\mathrm{full}}(y_t) - \log p_{\mathrm{alt}}(y_t)))$,
which couples the faithfulness requirement to gold-token correctness on
the full input. The model cannot satisfy the margin by being arbitrarily
wrong on the image-zeroed input.

The implementation runs three forward passes per training step --- full,
image-zeroed, and image-flipped --- and exposes a
\texttt{flipped\_row\_mask} that zeroes the flipped-counterfactual
contribution on rows that are not laterality-sensitive (claims that do
not contain one of \texttt{left}, \texttt{right}, \texttt{bilateral}).
Image-zeroing is implemented as the per-channel mean of the input batch's
pixel values, not a literal zero tensor: a literal zero is, after the
processor's normalisation, a specific out-of-distribution input the model
has never seen at training time; the per-channel mean produces a uniform
``neutral'' image at the model's own normalisation, which is what we want
when we say ``the image is absent.'' Edge cases: when no rows in the
batch are laterality-sensitive, the flipped term contributes zero; when
the assistant turn has zero target tokens, the loss is clamped to zero
without producing NaN gradients.

\subsection{Dual filter four-cell decomposition}
\label{app:filter}

\begin{table*}[t]
\centering
\small
\caption{\textbf{Four-cell decomposition of which filter branch flags each
training row.} On the $20{,}560$-row training pool, the image-only branch
catches an additional $6.5\%$ of rows that the text-only branch alone
would miss; this is the symmetric-extension contribution axis relative to
\citet{li2025sfc}.}
\label{tab:filter-decomp}
\begin{tabular}{lrr}
\toprule
\textbf{Filter outcome} & \textbf{Count} & \textbf{\% of pool}\\
\midrule
Both filters caught (unimodally solvable on either axis) & $15{,}375$ & $74.8\%$\\
Text-only caught, image-only passed                       & $\phantom{0}2{,}156$  & $10.5\%$\\
\textbf{Image-only caught, text-only passed (novelty axis)} & $\phantom{0}\mathbf{1{,}329}$  & $\mathbf{6.5}\%$\\
Neither caught (genuinely multimodal)                     & $\phantom{0}1{,}700$  & $\phantom{0}8.3\%$\\
\bottomrule
\end{tabular}
\end{table*}

After the subsequent text-target filter restricts training to OpenI rows
with non-empty findings text, the per-row weight distribution within the
surviving $10{,}839$ rows is $95.6\%$ at $w_{\min}=0.2$ and $4.4\%$ at
$w=1.0$. The continuous-SFC ablation replaces the discrete
$\{w_{\min}, 1\}$ rule with $\mathrm{SFC}(x) = 1 - p_{\text{text}}^{\text{correct}}(x)$
clipped to $[w_{\min}, 1]$, which provides a smoother weighting and
matches the published SFC recipe.

\subsection{IMG\_correct full derivation, edge cases, and the
``IMG\_correct $>$ IMG'' phenomenon}
\label{app:imgcorrect}

For a binary classifier $f$ with gold class $y$ on input $x$, define four
mutually-exclusive transitions under image-masking:
\begin{itemize}\itemsep -2pt
\item $\mathrm{cc}(x)$: correct on full, correct on masked.
\item $\mathrm{ci}(x)$: correct on full, incorrect on masked (the IMG sub-population).
\item $\mathrm{ic}(x)$: incorrect on full, correct on masked.
\item $\mathrm{ii}(x)$: incorrect on both.
\end{itemize}
Then $\mathrm{IMG} = (n_{\mathrm{ci}} - n_{\mathrm{ic}})/n \times 100$,
because IMG measures $\mathrm{acc}_{\mathrm{full}} - \mathrm{acc}_{\mathrm{masked}}
= ((n_{\mathrm{cc}} + n_{\mathrm{ci}}) - (n_{\mathrm{cc}} + n_{\mathrm{ic}}))/n
= (n_{\mathrm{ci}} - n_{\mathrm{ic}})/n$. IMG\_correct further partitions
$\mathrm{ci}$ into a margin-passing subset and a margin-failing subset:
\begin{align}
n_{\mathrm{ci,no-margin}}      &= |\{i : \mathrm{ci}(x_i)\}|\\
n_{\mathrm{ci,with-margin}}    &= |\{i : \mathrm{ci}(x_i) \wedge p_{\mathrm{full}}(y_i) - p_{\mathrm{masked}}(y_i) \geq \tau\}|
\end{align}
and defines:
\begin{align}
\mathrm{IMG\_correct}    &= n_{\mathrm{ci,with-margin}}/n \times 100\\
\mathrm{induced\ inv.}   &= (n_{\mathrm{ci,no-margin}} - n_{\mathrm{ci,with-margin}})/n \times 100.
\end{align}

\textbf{IMG\_correct can exceed IMG.} When the transition population
$\mathrm{ic}$ is non-trivial (model became correct under masking) and the
margin-failing-$\mathrm{ci}$ population is small, IMG\_correct is
$n_{\mathrm{ci,with-margin}}/n$ which can be larger than
$(n_{\mathrm{ci}} - n_{\mathrm{ic}})/n$. This is mathematically why
several configurations in our retroactive audit (Table~\ref{tab:audit-full},
\S\ref{app:audit}) show IMG\_correct $>$ IMG. It does \emph{not} indicate
negative induced inversion. For autoregressive generators we apply this
same logic per token on the gold target sequence, average over assistant-
turn token positions only, and require the causal-LM shift between logits
and target IDs (logits at position $t$ predict token $t+1$).

\subsection{Training hyperparameters and compute}
\label{app:hyper}

\begin{table*}[t]
\centering
\small
\caption{\textbf{Training hyperparameters across all five configurations.}
The five configurations differ only in (i) which weights JSONL is loaded
(if any), (ii) whether the faith loss is enabled, (iii) whether the
flipped counterfactual is enabled, and (iv) the resulting
$\lambda_{\mathrm{faith}}$. All other hyperparameters are identical.}
\label{tab:hyper}
\begin{tabular}{lc}
\toprule
\textbf{Hyperparameter} & \textbf{Value}\\
\midrule
Backbone                          & MedGemma-4B-IT (research-only licence)\\
LoRA rank $r$                     & $16$\\
LoRA alpha                        & $32$\\
LoRA targets                      & q\_proj, k\_proj, v\_proj, o\_proj\\
Trainable parameters              & $11.9$\,M ($0.28\%$ of $4.31$\,B base)\\
Optimizer                         & AdamW\\
Learning rate                     & $2\times 10^{-4}$\\
Weight decay                      & $0.0$ on LoRA targets, $0.01$ elsewhere\\
LR schedule                       & cosine to zero, $50$-step linear warm-up\\
Batch size                        & $2$\\
Gradient accumulation             & $\times 8$ (effective batch size $16$)\\
Max optimizer steps               & $2{,}500$\\
Mixed precision                   & bf16\\
Faith-loss margin (nats)          & $0.3$\\
Image masking                     & per-channel mean of batch pixel values\\
$\lambda_{\mathrm{faith}}$        & $0.3$ when enabled, $0$ otherwise\\
$w_{\min}$ (filter downweight)    & $0.2$\\
Filter activation threshold       & $0.7$ on $p_{\text{branch}}^{\text{correct}}$\\
Random seeds                      & $\{42, 1337, 9001\}$\\
GPU                               & NVIDIA H100 80\,GB (Modal)\\
Wallclock per fine-tune           & $1.5$--$4.0$\,h (faith-loss configs $\sim 3\times$ slower)\\
\bottomrule
\end{tabular}
\end{table*}

The faith-loss configurations are $\sim 3\times$ slower per step than
simple-loss configurations because each step runs three forward passes
(full, image-zeroed, image-flipped). Total Phase 4 fine-tuning compute is
$\approx \$200$ on Modal H100 80\,GB across the $15$ fine-tunes
($5$ configs $\times$ $3$ seeds). Total project envelope (training,
primary evals, retroactive audit, cross-counterfactual, cross-site) is
$\approx \$400$. All experiments are reproducible on a single H100
$80$\,GB instance with the released code.

\subsection{Test coverage}
\label{app:tests}

The released codebase contains $30$ unit tests (CPU, $\sim 1$\,s
wallclock) covering: (i) faith-loss behaviour --- zero-loss when full
overwhelms alt; positive-loss when equal; SFT-coupling-required scenario;
correctness term; padding ignored; weights additive; differentiability;
flipped-row-mask filtering; zero-valid-rows no-NaN
(\texttt{tests/test\_faith\_loss.py}, $9$ tests). (ii) IMG\_correct
behaviour --- perfect grounding; pure inversion caught at margin; mixed;
no-change; shape mismatch raises; generator-shift contract; v5\_3
replication scenario (\texttt{tests/test\_img\_correct.py}, $7$ tests).
(iii) Dual-filter behaviour --- min, product aggregations; mismatched
row-sets raises; unknown aggregation raises; disagreement-rate counting
(\texttt{tests/test\_dual\_filter.py}, $5$ tests). (iv) Data-loader
behaviour --- laterality regex; collate padding; three image variants on
laterality batch; flipping disabled; no-laterality rows; labels equal
input-ids; verifier-only filter drop; weight-remapping with interleaved
no-target rows (\texttt{tests/test\_data\_loader.py}, $9$ tests).

\section{Full IMG\_correct Retroactive Audit}
\label{app:audit}

\begin{table*}[t]
\centering
\small
\caption{\textbf{Full IMG\_correct retroactive audit, all 18 cells} (nine
verifier configurations + nine v5 ablations, each on $n=2{,}974$ verifier
rows). Acc full and Acc masked are accuracies under the full and image-
zeroed inputs, respectively.}
\label{tab:audit-full}
\setlength{\tabcolsep}{4.0pt}
\begin{tabular}{lrrrrrr}
\toprule
\textbf{Config} & \textbf{IMG\_corr} & \textbf{IMG} & \textbf{Inv} & \textbf{Inv/IMG} & \textbf{Acc full} & \textbf{Acc mask}\\
\midrule
\multicolumn{7}{l}{\emph{Headline configurations.}}\\
v5\_0\_base                   & 3.83  & 1.28  & 0.17 & 13.2\% & 92.30 & 91.02\\
v5\_1\_ground                 & 3.40  & 2.62  & 0.13 & \phantom{0}5.1\% & 91.22 & 88.60\\
v5\_2\_real                   & 64.66 & 61.84 & 0.91 & \phantom{0}1.5\% & 92.60 & 30.77\\
\textbf{v5\_3\_contrast}      & \textbf{71.05} & \textbf{68.12} & \textbf{0.74} & \textbf{1.1\%} & 93.01 & 24.88\\
v5\_4\_final                  & 68.43 & 63.55 & 0.40 & \phantom{0}0.6\% & 91.69 & 28.14\\
v6\_0\_3site                  & 74.51 & 70.21 & 0.07 & \phantom{0}0.1\% & 91.06 & 20.85\\
v6\_0\_loo\_no\_openi         & 38.20 & 34.20 & 0.00 & \phantom{0}0.0\% & 87.42 & 53.23\\
v6\_0\_loo\_no\_padchest      & 43.95 & 41.80 & 1.31 & \phantom{0}3.1\% & 92.30 & 50.50\\
v6\_0\_loo\_no\_chestx        & 18.59 & 14.53 & 0.20 & \phantom{0}1.4\% & 81.27 & 66.75\\
\addlinespace
\multicolumn{7}{l}{\emph{v5 ablations.}}\\
abl\_scale\_25                & 69.07 & 65.64 & 0.71 & \phantom{0}1.1\% & 92.70 & 27.07\\
abl\_scale\_50                & 70.51 & 67.59 & 0.81 & \phantom{0}1.2\% & 92.80 & 25.22\\
abl\_scale\_100               & 66.24 & 62.81 & 0.61 & \phantom{0}1.0\% & 92.80 & 29.99\\
abl\_hothresh\_70             & 73.84 & 69.97 & 0.57 & \phantom{0}0.8\% & 92.60 & 22.63\\
\textbf{abl\_hothresh\_80}    & \textbf{77.98} & \textbf{74.78} & 0.74 & \phantom{0}1.0\% & 92.77 & 17.99\\
abl\_no\_contrast             & 63.28 & 60.02 & 0.87 & \phantom{0}1.4\% & 92.77 & 32.75\\
abl\_no\_ground               & 72.53 & 69.64 & 0.50 & \phantom{0}0.7\% & 92.97 & 23.34\\
abl\_no\_uncert               & 73.40 & 70.54 & 0.98 & \phantom{0}1.4\% & 93.17 & 22.63\\
\textbf{abl\_no\_consist}     & \textbf{\phantom{0}3.60}  & \textbf{\phantom{0}2.79}  & 0.30 & \phantom{0}10.8\% & 88.80 & 86.01\\
\bottomrule
\end{tabular}
\end{table*}

\textbf{Distribution of induced-inversion contributions.} Across all 18
cells: induced inversion has an absolute range of
$[0.00, 1.31]$\,pp and a relative range of $[0.0\%, 13.2\%]$ of the
corresponding IMG. The cell with the largest \emph{relative} induced
inversion is $\mathrm{v5\_0\_base}$ at $13.2\%$, with the smallest
absolute IMG ($1.28$\,pp); fractionally large because any flipping at all
under masking is fractionally large at this scale. The cell with the
largest \emph{absolute} induced inversion is
$\mathrm{v6\_0\_loo\_no\_padchest}$ at $1.31$\,pp ($3.1\%$ of IMG). No
cell exceeds the pre-registered $\geq 30\%$ threshold; no cell is even
close.

\textbf{The \texttt{abl\_no\_consist} cell is the cleanest disconfirmation
of the pre-registered concern.} Removing the consistency loss collapses
IMG\_correct from $\sim 70$\,pp (nominal v5 family) to $3.60$\,pp ---
near baseline --- while \emph{simultaneously} reducing induced inversion
from $\sim 0.7$\,pp to $0.30$\,pp. The pre-registered concern was that
the consistency loss \emph{caused} induced inversion. The empirical
pattern is: with the consistency loss, IMG\_correct is high and induced
inversion is small; without it, both are small. The consistency loss
contributes the IMG\_correct, not the induced inversion.

\section{Cross-Counterfactual Detail}
\label{app:cross-cf}

\begin{table*}[t]
\centering
\small
\caption{\textbf{Cross-counterfactual eval --- full numbers including
masked accuracy.} Each row uses the same sft\_full seed-42 checkpoint;
only the counterfactual transformation differs. The masked accuracy
column makes visible that under the trained mean and noise
counterfactuals the model drives image-zeroed accuracy to near-zero
($1.51\%$ and $13.38\%$ respectively), while under crop and occlude the
model retains $71.83\%$ and $75.20\%$ accuracy --- close to the
unperturbed $75.48\%$. The faith loss has trained the model to drive
$p(y\mid \text{neutral image}) \to 0$ on the trained counterfactual; it
has not trained the model to do so on counterfactuals it has not seen.}
\label{tab:cross-cf-full}
\setlength{\tabcolsep}{3.0pt}
\begin{tabular}{lcccc}
\toprule
\textbf{Counterfactual} & \textbf{IMG} & \textbf{IMG\_corr} & \textbf{Acc full} & \textbf{Acc mask}\\
\midrule
\textbf{mean (TRAINED)} & 73.97 & 74.00 & 75.48 & \phantom{0}1.51\\
noise                   & 62.10 & 62.20 & 75.48 & 13.38\\
crop                    & \phantom{0}3.64 & \phantom{0}4.81 & 75.48 & 71.83\\
occlude                 & \phantom{0}0.28 & \phantom{0}1.11 & 75.48 & 75.20\\
\bottomrule
\end{tabular}
\end{table*}

\textbf{Baseline anchor.} On the matching crop counterfactual, sft\_only
seed-42 yields IMG $= 0.33$\,pp, IMG\_correct $= 2.80$\,pp at
$73.82\%$ accuracy, induced inversion $= 0.26$\,pp ($n=105{,}921$ tokens).
The cross-counterfactual sft\_full $-$ sft\_only IMG\_correct gap on
crop is $4.81 - 2.80 = 2.01$\,pp; the gap on the trained mean
counterfactual is $74.00 - 2.38 = 71.62$\,pp (single-seed comparison; the
3-seed comparison gives $72.58$\,pp). The order-of-magnitude difference
between trained-mean and crop is the operational meaning of the
$\S\ref{sec:cross-cf}$ partial-circularity finding.

\section{Cross-Site Detail}
\label{app:crosssite}

\textbf{Silver target construction.} The ChestX-Det10 test rows in the
standard test set have empty \texttt{evidence\_text} fields. We constructed
\texttt{groundbench\_test\_chestx\_det10\_silver.jsonl} ($1{,}122$ rows)
by injecting CheXagent-2-3B \citep{chen2024chexagent} generations as
\texttt{reference\_report}; the silver-target nature is reported
explicitly. Reviewers may push on this caveat. PadChest-GR cross-site
\citep{castro2024padchestgr} requires a fresh CheXagent generation pass
and is deferred. \textbf{Token-level analysis.} The cross-site
$54.65\%$ token-level accuracy is $20$\,pp lower than in-distribution
OpenI accuracy, reflecting (i) genuine cross-site difficulty (different
camera vendors, different radiologist phrasing) and (ii) the silver-
target nature (the model is being graded against another model's output,
not radiologist gold).

\section{Pre-Registered Hypotheses}
\label{app:prereg}

\textbf{(i) Full method reduces IMG\_correct on OpenI-test by $\geq 5$\,pp
vs SFT-only at matched accuracy.} Status: \textbf{HIT BY AN ORDER OF
MAGNITUDE} on three seeds $\times$ three faith-loss configs. Realised:
$+72.58$\,pp at $+1.17$\,pp accuracy.

\textbf{(ii) IMG\_correct refinement reveals $\geq 30\%$ of $\mathrm{v5\_3\_contrast}$'s
IMG as induced inversion.} Status: \textbf{REFUTED}. Realised: $1.1\%$.
Pattern holds across all 18 audited cells (induced-inversion contributions
$\le 1.31$\,pp absolute, $\le 13.2\%$ relative).

\textbf{(iii) Dual filter dominates the faith loss in the ablation.}
Status: \textbf{INVERTED}. The empirical pattern is the opposite ---
sft\_dual is statistically indistinguishable from sft\_only baseline; the
faith loss is the load-bearing component (sft\_faith alone gives
$+71.73$\,pp over baseline); the increment from adding the dual filter on
top of the faith loss is $+0.85$\,pp seed-mean (within one standard
deviation of sft\_full). We retain the dual filter in the recipe because
the four-cell decomposition (Appendix~\ref{app:filter}) demonstrates a
symmetric-extension contribution axis that is not addressable by text-
only filtering, and the full method's $+0.85$\,pp seed-mean over
sft\_faith is small but consistent across all three seeds.

\textbf{Reporting commitment.} One pre-registration HIT, one REFUTED,
one INVERTED --- all reported transparently. Each was tested
independently in \S\ref{sec:results}. The methodological refinement
(IMG\_correct) lands as a contribution independently of the empirical
results, since it is a reusable audit-tool that improves the
interpretation of any future or prior IMG-style metric.

\section{Compute Disclosure}
\label{app:compute}

Total compute spent on this paper: $\approx \$400$ on Modal H100 80\,GB
(approximately $30$ H100-hours total --- $15$ Phase 4 fine-tunes at
$\sim 1.5$\,h each on simple-loss configurations and $\sim 3$--$4$\,h
each on faith-loss configurations; $15$ Phase 5 IMG\_correct evals at
$5$--$10$\,min each; $18$ retroactive audit cells at $3$--$5$\,min each;
$4$ cross-counterfactual cells; $3$ cross-site cells). All experiments are
reproducible on a single H100 $80$\,GB instance with the released code;
the budget is well within reach for any researcher with comparable cloud-
compute access. The full project envelope (including image-build retries,
smoke tests, and the inevitable launcher overhead from two Modal billing-
cycle cap-hit cycles during the run) is approximately $\$250$ over
baseline; the headline experiments themselves are $\le \$150$.

\end{document}
